\section{Discusión}\label{sec:discusion}

Los resultados de la sección~\ref{sec:resultados} cuentan dos historias que
conviene separar. Una es espacial y es robusta: la siniestralidad registrada
en la ZMM se concentra en corredores, la estructura persiste entre años y
sobrevive a la escala y a la corrección por pruebas múltiples. La otra es
temporal y es administrativa: el nivel de cada municipio cambia de un año a
otro por razones que no tienen que ver con el tránsito.

\subsection{El régimen de reporte municipal como fuente principal de error}\label{sec:disc-regimen}

El supuesto implícito de cualquier modelo de conteos sobre registros
administrativos es que el subregistro existe pero es estable: si un municipio
captura la mitad de sus accidentes, los captura así todos los años, y el sesgo
se absorbe en el intercepto. La literatura vial documenta que el subregistro
es la norma y no la excepción, y que su nivel varía con la severidad y con el
país \parencite{elvik1999}. Lo que esta base muestra es que en la ZMM
tampoco es estable en el tiempo. Santa Catarina pasa de 3 registros en 2020 a
cerca de tres mil en 2023; Guadalupe alcanza su máximo en plena pandemia y cae
a un tercio en 2024; Salinas Victoria alterna entre dos niveles año con año.
Ninguna dinámica vial produce saltos de ese tamaño en un solo periodo: lo que
cambia es la práctica de captura.

Tres resultados independientes apuntan a la misma causa. Primero, el residual
de la Binomial Negativa tiene forma de municipio: su $I$ de Moran cae de 0.31
a 0.07 al quitar la media municipal (sección~\ref{sec:res-conteo}), de modo que
lo que parecía un proceso espacial fino que faltaba modelar eran 18 bloques
con un nivel distinto cada uno. Eso explica también por qué el rezago espacial
no aportó nada: las vecinas de una celda comparten municipio y se equivocan en
la misma dirección. Segundo, el exceso de ceros sobre la Binomial Negativa se
concentra en unos pocos pares municipio$\times$periodo, con Santa Catarina antes
de 2022 a la cabeza; son celdas cuyo municipio no reportaba, no celdas sin
riesgo. \textcite{lord2005} advierten que la inflación de ceros en datos de
accidentes suele ser un artefacto de subregistro o de una unidad de análisis
mal elegida y no un segundo proceso generador; aquí el diagnóstico es
explícito. Tercero, dos años recientes de historia predicen 2024 mejor que
cinco, lo que solo tiene sentido si los años viejos describen un régimen de
reporte que ya no rige.

La consecuencia de diseño fue convertir el diagnóstico en estructura: el
modelo jerárquico separa un efecto municipal dinámico $m_{j,a}$, que absorbe
\emph{cuánto se reporta}, de un efecto espacial $u_i$ por celda, que absorbe
\emph{dónde, dentro del municipio, se concentra}. La separación tiene un
precio que hay que declarar: $m_{j,a}$ no es siniestralidad. Es la mezcla
inseparable de riesgo real y de intensidad de captura, y con esta base no hay
forma de repartirla. Lo que sí se gana es que $u_i$ queda neto del régimen:
es el primer producto del proyecto que compara a Guadalupe con Guadalupe.

La misma lógica pone un techo al pronóstico. El régimen de 2024 no está en los
datos de 2019--2023, y ningún modelo entrenado con el pasado puede anticipar
que un municipio empiece o deje de reportar. El pronóstico del jerárquico es
por persistencia —toma el nivel municipal de 2023— y sus intervalos son, por
construcción, un piso de la incertidumbre real: no incluyen el riesgo de
cambio de régimen. La alternativa natural, propagar una innovación $t_3$ más,
no está disponible: $\E[e^{T}] = \infty$ para cualquier $t$ de Student, y la
media predictiva y el CRPS dejarían de existir. Que el jerárquico cubra el
87\,\% con un intervalo nominal del 90\,\% aun sin ese término indica que el
resto de la incertidumbre está bien representada; el faltante es precisamente
el régimen.

\subsection{Lo que la base puede y no puede decir}\label{sec:disc-alcance}

La distinción anterior permite delimitar con precisión el alcance de los
resultados.

\paragraph{Lo que puede decir.} Dónde se concentra el \emph{reporte} de
accidentes, con qué estructura espacial y con qué persistencia. El índice de
Gini, el $I$ de Moran y los mapas locales con FDR coinciden a 250, 500 y
1\,000\,m de lado, lo que descarta que la estructura sea un artefacto de la
rejilla \parencite{openshaw1984}. El ranking de celdas de 2019 correlaciona 0,83
con el de 2024, con una pandemia de por medio. Y, dentro de cada municipio,
la base sostiene un perfil de riesgo relativo $u_i$ con la incertidumbre que
corresponde a cada celda: las celdas con pocos accidentes no se encogen mucho
en su estimación puntual, pero su intervalo es tres veces más ancho que el de
las celdas con más de cien, y esa diferencia es información para decidir.

\paragraph{Lo que no puede decir.} Primero, no distingue un municipio donde
bajó el riesgo de uno que dejó de capturar accidentes; quien use $m_{j,a}$ o
las series municipales como indicador de seguridad vial tiene que resolver esa
ambigüedad por fuera, con otra fuente. Segundo, no anticipa cambios de régimen.
Tercero, no dice \emph{por qué} una celda es riesgosa: sin exposición (tránsito,
kilómetros de vialidad) un corredor con muchos accidentes puede ser un corredor
peligroso o un corredor transitado. La exposición es la
covariable central de los modelos de frecuencia de accidentes
\parencite{lord2010}; sin ella, lo que aquí se estima es frecuencia y no tasa.
Cuarto, no permite inferencia causal sobre intervenciones: todas las
covariables son estructurales y ninguna es manipulable.

\paragraph{El PAI como síntesis.} El índice de precisión predictiva al 5\,\%
vale 8.4 en todos los modelos, desde el conteo histórico hasta el jerárquico.
Es la expresión cuantitativa de los dos párrafos anteriores. El \emph{dónde} lo
decide el historial, que es lo que la base mide bien; las
mejoras del jerárquico están en el \emph{cuánto} y en el \emph{con qué
	certeza} (CRPS 6.17 contra 6.56), que es donde el régimen pesa. El PAI es la
métrica habitual para juzgar mapas de puntos calientes \parencite{chainey2008}, y
moverlo exige covariables que varíen por celda y que el historial no contenga, y esas
no están en ATUS.

\subsection{Limitaciones}\label{sec:disc-limitaciones}

\paragraph{Geocodificación por intersección.} ATUS georreferencia la
intersección reportada, no el punto físico del choque: el 97,6\,\% de las
coordenadas distintas corresponde a un único par de calles. La rejilla de
500\,m tolera bien ese error de posición, pero cualquier análisis a escala de
cuadra o de tramo hereda un dato discreto sobre unos 60\,000 nodos, no un
patrón de puntos continuo.

\paragraph{Cambio de procedimiento en 2022.} La matriz de Spearman entre años
tiene dos bloques, 2019--2021 y 2022--2024, que coinciden con la aparición del
campo \texttt{OID} en los archivos del INEGI. La hipótesis más económica es un
cambio en el procedimiento de georreferenciación. Se incorporó como indicador
de periodo, pero no se pudo verificar con documentación de la fuente; si la
causa fuera otra, la lectura de la discontinuidad cambiaría, aunque no la del
régimen municipal, que opera sobre el volumen y no sobre la posición.

\paragraph{Faltantes no aleatorios.} Los campos con códigos centinela tienen
un mecanismo MNAR: en accidentes fatales el aliento alcohólico se ignora 2.4
veces más seguido (sección~\ref{sec:datos-calidad}). El modelo de conteo no
usa esos campos, así que no lo afecta; cualquier modelo de severidad futuro
tendrá que tratarlos explícitamente.

\paragraph{Una sola ventana de validación.} La validación fuera de muestra se
hizo en un solo año, 2024. Es la única prueba honesta disponible —un
jerárquico con 2\,241 efectos de celda ajusta el pasado casi a la perfección—,
pero un solo año no permite estimar la variabilidad de las métricas entre
ventanas, y 2024 pudo ser un año con más o menos cambios de régimen que el
típico.

\paragraph{Escala anual.} El jerárquico se ajusta sobre celda$\times$año. El
supuesto de separabilidad espacio-tiempo está apoyado (el efecto de mes es
común a todas las celdas y aportó casi nada), pero el modelo no dice nada
sobre estacionalidad local ni sobre horarios, que son preguntas operativas.

\paragraph{Cobertura territorial.} Se retiraron las celdas sin accidentes en
6 años y las 103 celdas isla. La decisión se justifica porque una celda vacía
en 379 mil registros no tiene red vial, pero implica que el modelo no
predice para territorio donde la urbanización todavía no llega, que es
donde crece la periferia.

\subsection{Caminos que no funcionaron}\label{sec:disc-caminos}

Se documentan porque cada uno descarta una explicación y porque, sin el
registro, alguien volvería a intentarlos.

\paragraph{BYM2.} La especificación original del nivel de celda era BYM2
\parencite{riebler2016}, que separa un componente espacial ICAR de uno
independiente y hace interpretable la fracción espacial $\phi$. No se
identificó en ninguna de sus dos parametrizaciones: la no centrada llevó
$\phi$ a 0 y la centrada a 1, con $\hat{R}$ de 2.06 y 1.65. El motivo es de
estructura y no de ajuste fino del muestreador. Con 5 años de conteos por
celda los datos fijan cada $u_i$ casi por completo, y el reparto entre
``espacial'' e ``independiente'' solo lo sostienen los priors. La estructura
espacial existe —el Moran de la tasa por celda neta del municipio es 0.41—; lo
que falta es información para descomponerla. El CAR propio, con una sola
varianza y un solo $\alpha$, se identifica sin problema. La lección es general:
BYM2 está pensado para mapeo de enfermedades con conteos escasos por área, y
pierde su ventaja cuando cada área tiene información de sobra.

\paragraph{Rezago espacial en el GLM.} Incluir el promedio de las vecinas como
covariable, con un Moran de 0,65 en los conteos, parecía obvio. Su coeficiente
fue $-0{,}002$ ($t = -0{,}6$). Fue el primer indicio de que el residual no era
espacial fino sino municipal (sección~\ref{sec:disc-regimen}).

\paragraph{Inflación de ceros.} El ZINB con inflación constante mejora el AIC,
pero reparte la probabilidad de cero estructural entre todas las celdas por
igual, cuando el exceso está concentrado en unos pocos municipios y periodos.
Ganar AIC con un mecanismo equivocado no es mejorar el modelo; el fenómeno se
modeló donde ocurre, en el nivel municipal.

\paragraph{Innovaciones normales.} Con una caminata aleatoria normal para el
municipio, $\sigma_{rw}$ salió en 0.61: inflada por los pocos municipios que
saltan, y por tanto demasiado ancha para los que apenas se mueven. Con
innovaciones $t_3$ bajó a 0.29, que describe al municipio típico, y los saltos
quedan en las colas.

\paragraph{Deriva municipal.} Una tendencia local por municipio, pensada para
capturar el crecimiento de la periferia, se encogió a cero ($\sigma_d = 0{,}07$),
introdujo 46 divergencias y no mejoró el CRPS. Con cinco años y cambios de
régimen, los datos no distinguen una tendencia de una racha de innovaciones.

\paragraph{Ventanas de historia por AIC.} Cuando el GLM recibe la historia
completa y la reciente, elige la completa, porque en entrenamiento explica
mejor 2019--2021 al incluirlos; fuera de muestra pierde. Es la razón por la que
ninguna comparación del proyecto se hace con criterios de información sobre
el periodo de ajuste, sino con la validación en 2024.

\paragraph{Lo que no se intentó.} La función $K$ de Ripley y la densidad kernel
en su versión plana suponen que un accidente puede ocurrir en cualquier punto
del plano; aquí solo puede ocurrir sobre la red vial, y el nulo plano
sobreestimaría el agrupamiento. Las versiones restringidas a red
\parencite{okabe2001} exigen la geometría de las calles, que el proyecto no tiene.
Habrían confirmado algo ya establecido —que hay agrupamiento y a qué escalas—,
así que se pospusieron hasta tener la red vial, que es también la covariable
de exposición que falta (sección~\ref{sec:conclusiones}).
